\scanpage{30}
8) Soient $\ell,\ell'\in\mathbb P$ premiers aux car rés. % L1
de $X$. Soit $M\in\mathrm{Ob}\,\mathcal M^+(X)$. On veut % L2
que % L3
\[
\begin{gathered}
T_\ell(M)\ \text{faisceau constant tordu}\\ % L4
\Updownarrow\\ % L5
T_{\ell'}(M)\ \text{faisceau constant tordu} % L6
\end{gathered}
\]
et que sous ces conditions, $T_\ell(M)$ et % L7
$T_{\ell'}(M)$ doivent avoir $\unsure{\bar m|=}$ rang en chaque % L8
pt. / Pour vérifier l'égalité des rangs, on % L9
est ramené au cas où $X$ est \inserted{le} spectre d'un % L10
corps (fini si on veut, \deleted{\unsure{d}} ou clôture alg. % L11
d'un tel). Plus généralement, si % L12
$u$ est un endomorphisme de $M$, \deleted{\unsure{on|et}} % L13
($X=\mathrm{Spec}\,k$, $k$ un corps), on en déduit % L14
\[
T_\ell(u)\in\mathrm{End}\,T_\ell(M),\qquad T_{\ell'}(u)\in\mathrm{End}\,(T_{\ell'}(M)) % L15
\]
et je dis que on a % L16
\[
\boxed{\mathrm{Tr}\,T_\ell(u)=\mathrm{Tr}\,T_\ell(\unsure{u'|u})\in\deleted{\unsure{\mathbb Z}}\mathbb Q} % L17
\]
d'où, remplaçant $u$ par $\bigwedge^i u$, le fait % L18
\[
\boxed{P(T_\ell(u),t)=P(T_{\ell'}(u),t)\in\mathbb Q[t]} % L19
\]

% archive sheet number: 29
